\section{Objetivo y montaje}

La validación física se planteó para comprobar el funcionamiento del sistema
completo sobre una Basys 3. A diferencia de la simulación y del análisis
posterior al ruteo, el ensayo incluyó el cable Micro-USB, la interfaz JTAG, el
puente USB--UART, el oscilador de la placa y los LED utilizados como salida. El
alcance contempló el recorrido de una transacción desde la computadora hasta la
ALU y su retorno como respuesta UART, junto con la indicación visual del
resultado y las banderas.

Con la placa apagada, el selector de alimentación se colocó en USB y el cable
Micro-USB se conectó a J4, identificado como \texttt{PROG}. Luego se accionó
SW16 y se comprobó el encendido de LD20. El mismo enlace proporcionó
alimentación, acceso JTAG y comunicación USB--UART, tal como se describe en el
\href{https://digilent.com/reference/_media/reference/programmable-logic/basys-3/basys3_rm.pdf}{manual de referencia de la Basys 3}.

\section{Programación y configuración}

El proyecto \texttt{vivado/tp2\_uart/tp2\_uart.xpr} se abrió en Vivado y se
estableció la conexión desde \textit{Open Hardware Manager} mediante
\textit{Open Target $\rightarrow$ Auto Connect}. Hardware Manager detectó el
dispositivo \texttt{xc7a35t\_0}, sobre el cual se cargó
\texttt{build/alu\_uart\_top.bit} con \textit{Program Device}. Se adoptaron como
criterios de aceptación la finalización sin errores y el encendido permanente
del indicador \texttt{DONE}. El flujo coincidió con el procedimiento de
conexión y programación indicado en
\href{https://docs.amd.com/r/2025.1-English/ug908-vivado-programming-debugging/Opening-Hardware-Target-Connections}{Vivado UG908}.

\placeholder{\texttt{vivado\_hardware\_manager\_placa.png}}
{Dispositivo \texttt{xc7a35t\_0}, estado programado y archivo bitstream.}

Después de la programación se presionó y se soltó BTNC. El patrón de referencia
posterior al reset consistió en LED0--LED7 apagados, LED8 encendido y
LED9--LED10 apagados: resultado \texttt{00}, bandera Z activa y banderas C y V
inactivas.

El sistema operativo enumeró los puertos serie disponibles mediante:

\begin{codebox}
python -m serial.tools.list_ports -v
\end{codebox}

El canal asociado al puente USB--UART se asignó a la variable
\texttt{UART\_PORT}. El cliente utilizó \num{19200} baud, ocho bits de datos,
sin paridad, un bit de stop y sin control de flujo.

\section{Pruebas realizadas}

\subsection{Casos dirigidos}

La primera ejecución aplicó trece casos dirigidos que cubrieron las ocho
operaciones, resultados nulos, carry, overflow y opcode inválido. Cada respuesta
recibida se comparó automáticamente con el modelo de referencia y la salida
completa quedó registrada en \texttt{reports/placa\_dirigida.log}.

\begin{codebox}
python tools/uart_alu_client.py --port "$UART_PORT" \
  --suite directed 2>&1 | tee reports/placa_dirigida.log
\end{codebox}

Se estableció como criterio de aceptación que las trece comparaciones
coincidieran y que la ejecución terminara con:

\begin{codebox}
PASS: 13 operaciones completadas sin errores
\end{codebox}

\subsection{Casos pseudoaleatorios}

La segunda ejecución aplicó cien pares de operandos por cada una de las ocho
operaciones. La semilla \texttt{0x1a2b3c4d} fijó la secuencia y permitió
reproducir las 800 transacciones. La salida se conservó en
\texttt{reports/placa\_aleatoria.log}.

\begin{codebox}
python tools/uart_alu_client.py --port "$UART_PORT" \
  --suite random --random-per-op 100 --seed 0x1a2b3c4d \
  2>&1 | tee reports/placa_aleatoria.log
\end{codebox}

El criterio exigió que los 800 resultados recibidos coincidieran con los valores
esperados y que el cierre fuera:

\begin{codebox}
PASS: 800 operaciones completadas sin errores
\end{codebox}

\subsection{Recuperación después de timeout}

La recuperación se ejercitó mediante el envío de un único byte
\texttt{AA}. Después se esperaron \SI{20}{\milli\second}, un intervalo superior
al timeout interno de \SI{10}{\milli\second}, y se transmitió la resta
\texttt{09 - 02}.

\begin{codebox}
python tools/uart_alu_client.py --port "$UART_PORT" \
  --recovery 2>&1 | tee reports/placa_recuperacion.log
\end{codebox}

El criterio de aceptación exigió que el byte aislado no produjera una respuesta
y que la transacción completa posterior devolviera únicamente \texttt{07}. Ese
comportamiento indicaría que el controlador descartó el estado parcial antes de
recuperar la espera del primer operando.

\subsection{Correspondencia entre UART y LED}

La correspondencia entre la respuesta serie y la salida física se comprobó con
dos sumas. La primera operación produjo simultáneamente resultado nulo y carry:

\begin{codebox}
python tools/uart_alu_client.py --port "$UART_PORT" \
  --a 0xff --b 0x01 --opcode ADD \
  2>&1 | tee reports/placa_led_cz.log
\end{codebox}

El resultado esperado en la computadora fue \texttt{00}. El patrón físico de
referencia encendió LED8 y LED9 y mantuvo apagados LED0--LED7 y LED10, lo que
formó la palabra \texttt{0x300}.

\begin{figure}[H]
	\centering
	\fcolorbox{placeholderline}{placeholdergray}{%
		\parbox[c][5.3cm][c]{0.78\textwidth}{%
			\centering
			\textbf{Fotografía pendiente}\\[0.25cm]
			\texttt{basys3\_uart\_carry\_zero.jpg}\\[0.2cm]
			\small Placa completa; LED8 y LED9 encendidos.
		}%
	}
	\caption{Estado de la placa para la suma
	\texttt{FF + 01}: carry y zero activos.}%
\end{figure}

La segunda operación ejercitó el overflow de una suma con operandos positivos:

\begin{codebox}
python tools/uart_alu_client.py --port "$UART_PORT" \
  --a 0x7f --b 0x01 --opcode ADD \
  2>&1 | tee reports/placa_led_v.log
\end{codebox}

La respuesta UART esperada fue \texttt{80}. El patrón físico de referencia
encendió LED7 y LED10 y mantuvo apagados los demás LED, lo que correspondió a la
palabra \texttt{0x480}. Los dos casos permitieron cotejar el resultado
transmitido con el retenido en las salidas físicas.

\begin{figure}[H]
	\centering
	\fcolorbox{placeholderline}{placeholdergray}{%
		\parbox[c][5.3cm][c]{0.78\textwidth}{%
			\centering
			\textbf{Fotografía pendiente}\\[0.25cm]
			\texttt{basys3\_uart\_overflow.jpg}\\[0.2cm]
			\small Fila de LED; LED7 y LED10 encendidos.
		}%
	}
	\caption{Estado de los LED para la suma
	\texttt{7F + 01}: overflow activo.}%
\end{figure}

Al finalizar se reunieron las últimas líneas de los cinco registros para
conservar un resumen de la ejecución:

\begin{codebox}
clear
tail -n 1 reports/placa_dirigida.log
tail -n 1 reports/placa_aleatoria.log
tail -n 1 reports/placa_recuperacion.log
tail -n 1 reports/placa_led_cz.log
tail -n 1 reports/placa_led_v.log
\end{codebox}

\placeholder{\texttt{validacion\_uart\_terminal.png}}
{Resultado final de los cinco registros de ejecución.}

\section{Resultados y criterio de aceptación}

Para cada prueba se definió un resultado esperado antes de contrastar la salida
UART y los LED. Los resultados observados y su estado se mantuvieron pendientes
hasta disponer de los registros y las imágenes obtenidos sobre la placa.

\begin{table}[H]
	\centering
	\small
	\begin{tabularx}{0.96\textwidth}{L{2.7cm}Y C{2.4cm}C{1.9cm}}
		\toprule
		\textbf{Prueba} & \textbf{Criterio} &
		\textbf{Resultado observado} & \textbf{Estado} \\
		\midrule
		Reset & Resultado \texttt{00}; solo Z activa. & Pendiente & Pendiente \\
		Suite dirigida & 13 operaciones sin errores. & Pendiente & Pendiente \\
		Suite aleatoria & 800 operaciones sin errores. & Pendiente & Pendiente \\
		Recuperación & Una única respuesta \texttt{07}. & Pendiente & Pendiente \\
		ADD \texttt{FF+01} & UART \texttt{00}; LED = \texttt{0x300}. &
		Pendiente & Pendiente \\
		ADD \texttt{7F+01} & UART \texttt{80}; LED = \texttt{0x480}. &
		Pendiente & Pendiente \\
		\bottomrule
	\end{tabularx}
	\caption{Resultados de la validación física.}%
	\label{tab:validacion-fisica}
\end{table}

\section{Interpretación}

El conjunto de pruebas abarcó la cadena completa: enumeración del puerto,
recepción UART, almacenamiento temporal, interpretación de la terna, ejecución
de la ALU, transmisión de la respuesta y retención del resultado en los LED. La
suite dirigida ejercitó las situaciones de frontera y las 800 transacciones
pseudoaleatorias ampliaron el conjunto de operandos con una secuencia
reproducible.

La prueba de recuperación evaluó que una operación incompleta no contaminara la
transacción siguiente. Las dos comprobaciones visuales relacionaron la respuesta
recibida por la computadora con el estado de resultado, Z, C y V sobre la placa.
La aceptación final quedó determinada por la coincidencia de esas observaciones
con los criterios de la tabla~\ref{tab:validacion-fisica}.
